\documentclass[12pt,a4paper]{article}
\usepackage{amsmath}
\usepackage{amssymb}
\usepackage{bm}
\usepackage[utf-8]{inputenc}
\usepackage{xeCJK}

\setCJKmainfont{Hiragino Sans GB}

\title{京都大学大学院 数理工学専攻}
\author{凸最適化 過去問}
\date{H25 (OR)}

\begin{document}

\maketitle

\section*{解答 (H25)}

\subsection*{(i)}

$f$ と $g_j$ が凸関数であるため，$\ell(x)=f(x)+\sum_{j=1}^m\mu_j^*g_j(x)$ も凸関数である。
凸関数の性質により，任意の $x,y\in\mathbb{R}^n$ に対して
\[\ell(x) - \ell(y) \ge \nabla\ell(y)^T(x-y)\]
が成り立つ。

これは微分可能な凸関数の一階の充分条件である。
具体的には，$z:=y+t(x-y)$ とおき ($0<t<1$) ，凸性の定義
\[\ell(z) \le (1-t)\ell(y) + t\ell(x)\]
から
\[\frac{\ell(z)-\ell(y)}{t} \ge \ell(x) - \ell(y)\]
を得て，$t\to 0$ とすれば所望の不等式が従う。

\subsection*{(ii)}

KKT 条件より $\nabla\ell(x^*) = \nabla f(x^*) + \sum_{j=1}^m\mu_j^*\nabla g_j(x^*) = 0$ が成り立つ。

(i) の不等式に $y=x^*$ を代入すると
\[\ell(x) - \ell(x^*) \ge 0 \cdot (x-x^*) = 0\]
すなわち
\[\ell(x) \ge \ell(x^*)\]
が任意の $x\in\mathbb{R}^n$ に対して成り立つ。

\subsection*{(iii)}

任意の実行可能解 $x$ に対して $g_j(x)\le 0$ (すべての $j$) が成り立つ。
さらに補完性条件 $\mu_j^*g_j(x^*)=0$ より，右辺に現れる $\sum_{j=1}^m\mu_j^*g_j(x^*)=0$ である。

(ii) の結果を展開すると
\begin{align*}
f(x) + \sum_{j=1}^m\mu_j^*g_j(x) &\ge f(x^*) + \sum_{j=1}^m\mu_j^*g_j(x^*) \\
f(x) + \sum_{j=1}^m\mu_j^*g_j(x) &\ge f(x^*).
\end{align*}

$x$ が実行可能ならば $g_j(x)\le 0$ で，$\mu_j^*\ge 0$ であるから
\[\sum_{j=1}^m\mu_j^*g_j(x) \le 0.\]

したがって
\[f(x) \ge f(x^*) - \sum_{j=1}^m\mu_j^*g_j(x) \ge f(x^*)\]
が成り立ち，$x^*$ は問題 P の大域的最適解である。

\subsection*{(iv)}

KKT 条件は
\[
\begin{pmatrix} x_1-1 \\ x_2-1 \end{pmatrix} + \mu_1^* \begin{pmatrix} -1 \\ 0 \end{pmatrix} + \mu_2^* \begin{pmatrix} 1 \\ 1 \end{pmatrix} = \begin{pmatrix} 0 \\ 0 \end{pmatrix}
\]
すなわち
\begin{align*}
x_1-1 - \mu_1^* + \mu_2^* &= 0, \\
x_2-1 + \mu_2^* &= 0,
\end{align*}
ならびに
\[
-x_1 \le 0, \quad x_1+x_2-1 \le 0, \quad \mu_1^*\ge 0, \quad \mu_2^*\ge 0, \quad \mu_1^*(-x_1)=0, \quad \mu_2^*(x_1+x_2-1)=0.
\]

相補性条件を考察する。

\paragraph{場合 (I): $\mu_1^*=0$}
\begin{enumerate}
\item[(I-a)] $\mu_1^*=0$ かつ $\mu_2^*=0$ のとき：
  第1式から $x_1=1$，第2式から $x_2=1$ だが，$x_1+x_2-1=1\not\le 0$ で実行不可能。
  
\item[(I-b)] $\mu_1^*=0$ かつ $x_1+x_2-1=0$ のとき：
  第1式から $\mu_2^*=1-x_1$，第2式から $\mu_2^*=1-x_2$ である。
  両者が等しいので $x_1=x_2$。$x_1+x_2-1=0$ より $x_1=x_2=1/2$。
  このとき $\mu_2^*=1/2>0$，$x_1=1/2>0$ より $-x_1<0$ も満たす。
  したがって $x^*=(\tfrac{1}{2},\tfrac{1}{2})^T$，$\mu^*=(0,\tfrac{1}{2})^T$ が解となる。
\end{enumerate}

\paragraph{場合 (II): $-x_1=0$ (すなわち $x_1=0$)}
\begin{enumerate}
\item[(II-a)] $x_1=0$ かつ $\mu_2^*=0$ のとき：
  第2式から $x_2=1$，第1式から $-1-\mu_1^*=0$ より $\mu_1^*=-1<0$，条件違反。
  
\item[(II-b)] $x_1=0$ かつ $x_1+x_2-1=0$ のとき：
  $x_2=1$。第1式から $-1-\mu_1^*+\mu_2^*=0$，第2式から $\mu_2^*=0$ より
  $\mu_1^*=-1<0$，条件違反。
\end{enumerate}

\paragraph{結論}

唯一の解は
\[
x^* = \begin{pmatrix} 1/2 \\ 1/2 \end{pmatrix}, \qquad \mu^* = \begin{pmatrix} 0 \\ 1/2 \end{pmatrix}.\]

\end{document}
